\mysection{ループ}

プログラムが何らかのファイルや特定のサイズのバッファを扱う場合、
コード内には何らかの復号化/処理ループが必要です。

これは\tracerツールの出力の実際の例です。
258バイトの暗号化されたファイルをロードするコードがありました。
各命令の実行回数を取得する目的で実行しました（最近では\ac{DBI}ツールの方がはるかに良いサービスを提供します）。
そして、259/258回実行されるコード片をすぐに見つけました：

\lstinputlisting{digging_into_code/crypto_loop.txt}

これが復号化ループであることがわかりました。